
\section{Methodology and Data}

Section outline

\subsection{Rubin Data Preview 2} 

Intro to DP2

We use the COSMOS field because of the availability of external
COSMOS2020/FARMER labels \citep{2022ApJS..258...11W}. For the COSMOS
validation sample, we select Rubin DP2 sources in the rectangular region
\(149.7 < {\rm RA} < 150.7\) and \(1.7 < {\rm Dec} < 2.7\), with
\(16 \le r_{\rm CModel} < 26\). External star--galaxy labels are taken from
the COSMOS2020/FARMER reference catalog and matched to DP2 sources using a
0.5 arcsec matching radius. After applying the same region and magnitude cuts,
the DP2 sample contains 269,606 sources, of which 193,586 are matched to
external labels and 76,020 are unmatched. The raw DP2-side matching fraction is
therefore \(193{,}586/269{,}606 \simeq 0.718\). The matched validation sample
contains 8,788 stars and 184,798 galaxies, using \(\texttt{truth\_binary}=1\)
for stars and \(\texttt{truth\_binary}=0\) for galaxies.

\begin{figure*}[t]
  \centering
  \includegraphics[width=0.5\textwidth]{fig01_cosmos_sky_coverage.png}
  \caption{
  COSMOS validation footprint. Gray points show selected Rubin DP2 sources in
  the common COSMOS region, and green points show the subset matched to
  COSMOS2020/FARMER external labels using a 0.5 arcsec matching radius.
  }
  \label{fig:cosmos_sky_coverage}
\end{figure*}

\begin{figure*}[t]
  \centering
  \includegraphics[width=0.5\textwidth]{fig02_cosmos_dataset_overview.png}
  \caption{
  Magnitude distribution of the COSMOS validation sample. The black curve shows
  all selected DP2 sources, while the blue and red curves show matched
  COSMOS2020/FARMER stars and galaxies. Counts are shown per magnitude using
  0.5 mag bins.
  }
  \label{fig:cosmos_counts}
\end{figure*}




\subsection{The distributions of sources classified using HST labels}

The external COSMOS2020/FARMER labels provide the reference classification used
to examine the morphology and color distributions of stars and galaxies. These
distributions motivate the probabilistic star--galaxy separation model
developed below.

The top left panel in Figure~\ref{fig:truth_morphology_color} shows the distribution of stars and galaxies classified
using HST labels in the \(\Delta m = m_{\mathrm{PSF}} - m_{\mathrm{CModel}}\) vs. \(m_{\mathrm{CModel}}\)
diagram. Modeling this distribution is the key ingredient of the star-galaxy separation method developed here. 


\begin{figure*}
  \centering
  \includegraphics[width=\textwidth]{fig03_truth_morphology_color.png}
  \caption{COSMOS/COSMOS2020 truth-label morphology and color overview. The r-band CModel magnitude
    is uncorrected for ISM dust extinction, while the CModel colors are dust-corrected where applicable. XXX: almost perfect!
    however, the counts plot shows the same numbers of stars and galaxies at about r = 19 (which is roughly
    correect) but the top right panel doesn't seem to be consistent with that - galaxies begin to
    dominate only at r of about 22-23. WE NEED TO DOUBLE CHECK THIS! Also, remove the title (here
    and elsewhere - the journal editors don't like it.}
  \label{fig:truth_morphology_color}
\end{figure*}

The top right panel in Figure~\ref{fig:truth_morphology_color} shows that galaxies, in addition to
outnumbering stars at faint magnitudes, have on average bluer colors than stars. 

Color-color diagrams shown in the bottom two panels in Figure~\ref{fig:truth_morphology_color}
demonstrate that colors can be used to enhance morphological star-galaxy separation. 
Nevertheless, Figure~\ref{fig:truth_color_color} and the top left panel in
Figure~\ref{fig:truth_morphology_color} show that increased noise close to the faint
magnitude limit will cause deteriorated performance for both morphological and color-based
star-galaxy separation methods. 


\begin{figure*}
  \centering
  \includegraphics[width=\textwidth]{fig04_truth_color_color.png}
  \caption{Similar to Figure~\ref{fig:truth_morphology_color}, except that data are 
    split by r-band CModel magnitude
   to show deterioration close to the faint magnitude limit (left: 16 $<$ rmag $<$ 25; right: 25 $<$
   rmag $<$ 26).}
  \label{fig:truth_color_color}
\end{figure*}


\subsection{The distributions of sources classified using \texttt{extendedness} parameter} 

Difficulties with star-galaxy separation encountered at the faint end are illustrated in
Figure~\ref{fig:r_extendedness_confusion_cmd}, which compares subsamples classified
using HST ground-truth labels and \texttt{extendedness} parameter in the r band.
As evident, a significant fraction of faint sources within 2 magnitudes from the faint limit are 
misclassified when using default \texttt{extendedness} parameter. 

\begin{figure*}
  \centering
  \includegraphics[width=\textwidth]{fig06_r_extendedness_confusion_cmd.png}
  \caption{Performance of default \texttt{extendedness} parameter in the r band
    relative to HST ground-truth labels. The top left and bottom right panels show
    correctly classified unresolved and resolved sources in the r vs. g-i color-magnitude
    diagram, with counts shown in the title. The bottom left and the top right panels
    show misclassified sources. Note that the stellar sample selected by \texttt{extendedness}
    is severely contaminated by galaxies for \(m_{\mathrm{CModel}} > 24\).}
  \label{fig:r_extendedness_confusion_cmd}
\end{figure*}




\subsection{Bayesian Star-Galaxy Separation Method}

Summary of the method, expressions from SlL2020... 


\subsection{Implementation of the SIL2020 method}

%Figure~\ref{} 


\subsubsection{Mixture Model for \(Delta m\) vs. \(m_{\mathrm{CModel}}\) diagram} 

XXX We need Nirav to describe in detail how the fits are done (what bins, what weights, etc).
We also show figures that demonstrate performance (e.g. histograms and 2D residuals). 


\subsubsection{Priors for Star-Galaxy Counts ratio} 

Here we discuss how the normalization parameters from the best-fit mixture can be used
to set priors for the star-to-galaxy number ratio. We also compare to the HST curve, as well
as to TRILEGAL-based curve. 





